
% Introduction
% =============================
% This section introduces the context, motivation, and research questions for the article.
% TODO: Refine the research questions and add more recent references.
% TODO: Add a conceptual figure of the remote fuzzy assessment model (see /figuras/sistema_difuso.png)
% TODO: Suggest a table comparing traditional and fuzzy-based assessment (see Table~%\ref{tab:comparacion})
% TODO: Include a simple equation for fuzzy inference (see Eq.~%\ref{eq:fuzzy})

\section{Introduction}

The rapid transition to remote and hybrid learning environments, accelerated by global events such as the COVID-19 pandemic, has exposed significant challenges in the assessment of student learning. Traditional evaluation methods, often designed for face-to-face settings, struggle to capture the complexity and diversity of student engagement and performance in virtual contexts. This situation has intensified the need for innovative assessment models that are both fair and supportive of meaningful learning \cite{Lopez2021, Black2020}.

Conventional grading systems tend to emphasize final outcomes, such as exam scores, rather than the learning process itself. This focus on grades can undermine student motivation, encourage superficial learning strategies, and fail to recognize important soft skills like participation, self-assessment, and collaboration \cite{Nguyen2020}. In remote settings, these limitations are further exacerbated by reduced direct interaction and increased variability in student circumstances.

A learning-centered approach to assessment prioritizes the development of competencies and the continuous improvement of students, rather than merely assigning a numerical grade. However, implementing such an approach requires tools capable of handling the inherent subjectivity and uncertainty in evaluating soft variables. Fuzzy logic offers a promising solution by providing a mathematical framework to formalize expert judgment and automate the grading process without losing the nuance of human evaluation \cite{Santos2022, Wang2023}.

This article proposes and validates a continuous assessment model based on fuzzy logic, designed to support meaningful learning in remote or hybrid environments. The model integrates multiple variables—including participation, assignment submission, and self-assessment—into a fair and transparent grading system. The main objective is to demonstrate that fuzzy logic can bridge the gap between teacher judgment and automated evaluation, transforming grading into a tool for real learning support.

% Example of a simple fuzzy inference equation
\begin{equation}
	ext{Grade} = \frac{\sum_{i=1}^n w_i \cdot \mu_i(x)}{\sum_{i=1}^n w_i}
\label{eq:fuzzy}
\end{equation}

% Example of a table suggestion
% Table~%\ref{tab:comparacion} could compare traditional vs. fuzzy-based assessment criteria.

Conventional grading systems tend to emphasize final outcomes, such as exam scores, rather than the learning process itself. This focus on grades can undermine student motivation, encourage superficial learning strategies, and fail to recognize important soft skills like participation, self-assessment, and collaboration. In remote settings, these limitations are further exacerbated by reduced direct interaction and increased variability in student circumstances.

A learning-centered approach to assessment prioritizes the development of competencies and the continuous improvement of students, rather than merely assigning a numerical grade. However, implementing such an approach requires tools capable of handling the inherent subjectivity and uncertainty in evaluating soft variables. Fuzzy logic offers a promising solution by providing a mathematical framework to formalize expert judgment and automate the grading process without losing the nuance of human evaluation.


This article proposes and validates a continuous assessment model based on fuzzy logic, designed to support meaningful learning in remote or hybrid environments. The model integrates multiple variables—including participation, assignment submission, and self-assessment—into a fair and transparent grading system. The main objective is to demonstrate that fuzzy logic can bridge the gap between teacher judgment and automated evaluation, transforming grading into a tool for real learning support.\cite{dummy}

% TODO: Add a figure illustrating the conceptual model in /figuras
% TODO: Expand on the research questions and expected contributions
% TODO: Cite foundational works on fuzzy logic in education